%\documentclass[letter,12pt,english, draft]{article}
\documentclass[letter, 12pt, english]{article}
\usepackage{mikkel}
\usepackage[ruled,vlined]{algorithm2e}
\title{Experimental Verification of the Lonely Runner Conjecture}
\author{Mikkel Kj�r Jensen}
\date{\today}
\begin{document}
\pagestyle{headings}
\maketitle
\abstract{
The Lonely Runner conjecture states that given n runners, with pairwise  speeds, there will be a time where every runner will be at least $\frac{1}{n + 1}$ distance away from their common start position. The conjecture has been proved up to n = 6\cite{serra_thelonely}. Since the Lonely Runner conjecture has applications in both Graph theory, number theory, and View-Obstruction problems, it would be nice if it could be proved for all n. In this project I have tried to experimentally verify the conjecture, up to 1000 runners (for speeds under 1000) \todo{When I have results fill this out}
}
\listoffixmes
\tableofcontents
\section{Introduction}
\label{introduction}
The Lonely Runner conjecture was first stated by J. M. Wills in \cite{JMWills}, where it was presented as a purely formal problem in the field of Diophantine approximation. It remained so until 1998, where it was given its more memorable name, and intuitive interpretation, from \cite{Bienia97flows.view-obstructions}. As of 2010 the Lonely Runner conjecture has only been proved for up to $n$ = 6 \cite{serra_thelonely}.

There exist several equivalent interpretations of the Lonely Runner conjecture, and I will here state widely used versions\footnote{See \cite{ANote}}:
\begin{enumerate}
\item Imagine $n$ runners at a circular, unit length track\footnote{i.e. the track has a circumference of 1}, where every runner runs with constant, pairwise different speeds\footnote{These speeds can, without loss of generality be assumed to be in $\N$ - see \cite{Bienia97flows.view-obstructions} and Section \ref{integerSpeeds}}, there will then exist a time t where all runners will be at least $\frac{1}{n + 1}$ units away from their common start-point.\\

\item Imagine we have $n$ + 1 runners, then the conjecture states that for every runner there exists a time where all other runners are at least $\frac{1}{n + 1}$ units away from it. Everything else is the same as in the first interpretation \cite{Bienia97flows.view-obstructions}.\\
\end{enumerate}

More formally the problem can be stated thus: 
Given any $n$ positive integers $w_1, w_2, \ldots, w_n$, there exists a real number $x$ such that 
\eqn{
\label{eqa:lonelyRunner} \Vert w_i x\Vert \geq \frac{1}{n+1}
}

for $i = 1, 2, \ldots, n$, where for $y \in \R$, $\Vert y \Vert$ is the distance from $y$ to the nearest integer. I owe the above formulations of the conjecture to \cite{ANote}. In this report it is useful to think of the track as the one-dimensional Torus $\mathbb{T} = [0,1)$.

\subsection{Expectations of the Reader}
I expect the reader to have the mathematical knowledge required to understand both the formal and informal description of the Lonely Runner Conjecture, as well having an understanding of university level Computational Geometry\footnote{For instance having read \cite{citeulike:3347056} or similar}.

\subsection{Scope and Limitations}
\begin{itemize}
\item In this project I will not attempt to prove the Lonely Runner conjecture for all n. 
\item I will create a program that can verify whether a given configuration of runners is true for equation \eqaref{eqa:lonelyRunner}, and make it return a result, within a reasonable time, for up to 1000 runners. The final program will work with more than 1000 runners, but I will prioritise readability, maintainability and portability of the program instead of extreme optimisations. 
\item The program will include a system to test a range of runner configurations against equation \eqaref{eqa:lonelyRunner}.
\item I will only focus on runners with integer speeds - see \ref{integerSpeeds} on p. \pageref{integerSpeeds} for an in depth justification for this.
\end{itemize}

\subsection{Computer Specifications}
\label{specs}
In this section I define the computer requirements for running my program. As mentioned above, the Lonely Runner conjecture has direct applications for View-Obstruction problems and colouring of regular chromatic graphs. Researchers working with in these fields have no special reason to work with exceedingly powerful computers, or computers with special hardware.
 
Thus, any algorithm I devise and implement should be able to run on a standard computer with no special hardware\footnote{Say hardware dedicated to fast floating point calculations}. I will therefore design my verification program to run on a 2 GHz computer with the recommended amount of RAM for the computer's OS. 

I reserve the right to use third-party software libraries - as long as these can be freely distributed and installed.

\subsection{Terminology}
\label{Termonolgy}
\begin{figure}[H]
  \centering
  \includegraphics[width=0.3\textwidth, angle=90]{./images/circleZoneEPS}
  \caption{\label{circleZoneImg}An illustration of the Runner track. The greyed out part is the Zone}
\end{figure}

In this report I will refer to runners as male for simplicity.

\begin{defi}
\label{def:theZone}
In this report I will refer to the track interval [$\frac{1}{n + 1}$, $\frac{n}{n+1}$] as the Zone, where $n$ is the number of runners, and a runner who is in this interval, as ``being in the Zone''. see figure \ref{circleZoneImg} for an illustration of the Zone.
\end{defi}

\begin{defi}
\label{def:config}
Configurations of runners, will in this report refer to a set S of $n \in \N$ runners, where for all runners $r, r\prime \in S$, $r \neq r\prime$, $r_{speed}, r\prime_{speed} \in \N$, and $r_{speed} \neq r\prime_{speed}$.
\end{defi}

\hide{
\begin{defi}
\label{def:fakeRunner}
When talking about the fake runner, I will be referring to a runner that is not actually in race, but rather a runner with speed 1, who I have introduced in order to have a convenient termination criterion for the Geometrical algorithm\footnote{\pref{compGeo}}. \pref{termination} for more.
\end{defi}  
}
\subsection{Canonical Example}
Through the rest of this report I will use the following configuration
of runners to give an example:
We have 4 runners with the speeds of 3, 5, 6, 8

\subsection{Background Material Used}
\label{background}
For this project I have read the following papers dealing with the Lonely Runner conjecture: \emph{``View-Obstruction Problems''}\cite{Bienia97flows.view-obstructions}, \emph{``The lonely runner problem with seven runners''} \cite{serra_thelonely}, \emph{``Regular chromatic number and the lonely runner problem''} \cite{Barajas2007479}, \emph{``View-Obstruction Problems''} \cite{springerlink:10.1007/BF01832623}, \emph{``Tight Instances of the Lonely Runner''} \cite{Goddyn96tightinstances}, \emph{``A Note on the Lonely Runner Conjecture''} \cite{ANote} and \emph{``Invisible runners in finite fields''} \cite{invis}.

I have also used the follow support literature:
\emph{``Computational Geometry''} \cite{citeulike:3347056}, \emph{``Uniform Distribution of Sequences''} \cite{uniform} and \emph{``Kalkulus''} \cite{kalkulus}.

\subsection{Source Code and Test Data}
The source code and the data used in this report (in machine readable form) can be found in my GitHub repository at \underline{https://github.com/kjmikkel/LonelyRunner}.

\subsection{Overview}
In this section I give a quick overview of what the different sections in report will cover.
\begin{description}
\item[Section \ref{introduction}] The introduction to the report, which will introduce the Lonely Runner conjecture and explain my goals
\item[Section \ref{choiceOfMethod}] A discussion of the two different approaches that can decide whether or not equation \eqaref{eqa:lonelyRunner} holds for a given configuration
\item[Section \ref{implementation}] Implementation details and an analysis of the scalability of the implementation
\item[Section \ref{section:program}] Description of the program I have created and how to use it.
\item[Section \ref{test}] Description of the tests I have performed.
\item[Section \ref{results}] Interpretation of the results of the tests.
\item[Section \ref{conclusion}] The conclusion of the report, where I will sum up my findings and discuss future work.
\end{description}

\section{Possible methods}
\label{choiceOfMethod}

\subsection{Introduction}
I will in this subsection describe the different methods I have thought of using in this project, and discuss their merits and flaws.

\begin{description}
\item[Computational Geometry:] The intuitive description of the Lonely Runner conjecture (descibed in the Introduction) lends itself well to a geometrical interpretation, and leads me to wonder whether an algorithm based on this fact would be efficient.

\item[Number theoretic:] The Lonely Runner conjecture is, in its original formulation, a problem from number theory. It is therefore not reasonable to assume that a number theoretic approach could lead to great speedups in calculating a possible solution.
\end{description}

\subsection{Computational Geometry}
The following solution is based on the first intiutive description of the conjecture. It is clear that we are interested in the time interval when the runners are in the \zone. However, we are interested in the time they are in the \zone - i.e. when a runner with a speed $s$ enters and leaves the \zone.\\ 

Let $s$ be the speed and $t$ be the time, then from \ref{mainProb} we have 

\begin{equation}
\label{speedOne}
\begin{split}
s * t &= \frac{1}{n+1} \\
t &= \frac{1}{s * (n+1)}
\end{split}
\end{equation}

and 

\begin{equation}
\label{speedTwo}
\begin{split}
s * t &= \frac{n}{n+1} \\
t &= \frac{n}{s * (n+1)}
\end{split}
\end{equation}

From \ref{speedOne} and \ref{speedOne} it is clear that the time interval for the first the runner is in the \zone will be 
\begin{displaymath}
\left[\frac{1}{s * (n+1)}, \frac{n}{s * (n+1)}\right]
\end{displaymath}

More generally, a given runner with speed $s$ will be in the \zone at time 

\begin{displaymath}
\left[\frac{1 + k * (n+1)}{s * (n+1)}, \frac{n + k * (n+1)}{s * (n+1)}\right] 
\end{displaymath}

So now it becomes a matter find out whether there is a time where the time intervals for all the runners overlap. To do this I propose using a horizontal Plane Sweep algorithm. Since the Lonely Runner Conjecture has not been proven, we have to make sure that the plane sweep algorithm will always terminate, but always find a valid solution to the conjecture, if its exists.

My solution to this problem would be to terminate the algorithm if all the runners crosses the start-line at the same time - since everything after that will only be a repitition of what had already happened.

\subsubsection{Points}

In order to distingush these 3 cases we need to introduce 3 different points:
\begin{description}
\item[\comStart] The first point where the runner is $\frac{1}{n + 1}$ units away from the start line. In the following algorithm the \comStart will contain the time it took to reach that location, and which runner it belongs to.
\item[\comEnd] The first point after the \comStart where the runner is $\frac{1}{n + 1}$ units away from the start line. In the following algorithm \comEnd has the responcibility of finding the next interval its runner is at least $\frac{1}{n+1}$ units away from the startline. I place this responcibility on \comEnd, on the grounds it is the first point where we know the current interval is not going to give a solution to the Lonely Runner conjecture. \comEnd needs to know the time its runner will pass it, the amount of time used to pass $\frac{1}{n+1}$ units, and to which runner it belongs to.
\item[\comFin] The time at which the runner passes the startline. As explainted above \todo{explain it above}, this is needed in order to ensure that the algorithm will always terminate, even if there is no solution to the Lonely Runner in that instance.
\end{description}

\subsubsection{Order of points}
In order to ensure this approch will work, I introduce the order $\prec$. For the timepoints $p$ and $q$, $p \prec q$ iff \\

\begin{center}
$p_{time} < q_{time}$\\
or \\
$p_{time} == q_{time}$ and type($p$) == \comFin and type($q$) != type($p$)\\
or \\
$p_{time} == q_{time}$ and type($p$) == \comStart and type($q$) == \comEnd \\
or \\
$p_{time} == q_{time}$ and type($p$) == type($q$) and $p_{runner} < q_{runner}$
\end{center}

Using this order there is no ambiguity where to place the points, and we also avoid the degernate case where we have (n-1) runners are in the \zone, and the next two points are a \comEnd and a \comStart, both at the same time. Since the \zone is a closed set, it is clear that this instance should return a valid solution to the conjecture, but if \comEnd comes first, then the case would not be reported - therefore \comStart must be placed before \comEnd.

\subsubsection{Algorithm}
\begin{algorithm}[H]
\caption{MakeTimePoints}
\highlights
\SetKwData{start}{startTime}
\Input{A start-time \ti and the \unit used to run $\frac{1}{(\n + 1)}$ part of the track, the number \run of the runner, and the time queue \li}
\Output{The time queue \li, with a new \startT, \eT and \finish inserted for \run runner}
 
Make new \startT $start$ from \start as its start-time, and \run as its runner
  
Make new \eT $end$ from \start + \unit * \n as its start-time, \unit as its speed, and \run as its runner
  
Make new \finish $finish$ with \start + \unit * (n+1) as its time, and \run as its runner
  
Add $start$, $end$ and $finish$ to \li

\return \li
\end{algorithm}

\begin{algorithm}[H]
  \caption{FindLonelyRunnerTime}
  \highlights
  \SetKwData{rTime}{finishTime}
  \SetKwData{lFinish}{LastRunnerToFinish}
  \SetKwData{and}{and}
  \Input{A list \s which contains n pairwise different speeds for n runners}
  \Output{A time \ti where all n runners are at least $\frac{1}{n + 1}$ units away from the starting line, or a \no, indicating that no such time exists}
  
  Create and initilize \li
  
\End $\gets 0$

\inter $\gets 0$

\rTime $\gets -1$

\lFinish $\gets -1$

\n $\gets$ size(\s)

$runnerNum \gets 1$

\ForEach{speed $\in$ \s}{
  
  \unit $\gets \frac{1}{speed * (\n + 1)}$  
  
  \li $\gets$ MakeTimePoints(\unit, \unit, $runnerNum$, \li)
  
  $runnerNum += 1$
}

\While{\li is not empty}{
  \p $\gets$ firstPoint(\li)
  
  \If{\p is \startT}{
      \End = 0
      
      \inter $+= 1$
      
      \If{\inter == \n}{
        \return $\p_{time}$
      }
  }
  
  \If{\p is \eT}{
    \End = 0
    
    \inter $-= 1$
    
    \ti $\gets \frac{2}{\p_{\unit} * (\n + 1)} + \p.\ti$
    
    \li $\gets$ MakeTimePoints(\ti, $p_{\unit}$, $p_{\run}$, \li)
  }
  
  \If{\p is \finish}{
    \If{\rTime == $-1$}{
      
      \rTime $\gets \p_{time}$
      
      \lFinish $\gets 0$ 
      
    }
    
    \If{$\p_{time}$ == \rTime \and $\p_{\run}$ == \lFinish}{
   
     \End $+= 1$\;

     \lFinish $+= 1$
    
     
     \If{\End == \n}{  
       \return \no
     }
     
     
    }
    \uElse(hello){
      
      \rTime $\gets -1$
   
      \End $\gets 0$
      
      \lFinish $\gets -1$
      
    }
    
  }
}
\end{algorithm}


\subsection{Number Theoretic}


\section{Method used}
\label{method}

\subsection{Introduction}

\subsection{The method}

\subsection{Good-Turing smoothing}

\section{Implementation}
\label{implementation}

In this section I will try to describe the implementation choices I have made in this project.


\subsection{Choice of programming language}
I will in this section discuss the factors that decided my choice of programming languages for this project, as well as the final choice.

I have used the following parameters to decide my choice of programming language:
\begin{description}
\item[Speed:] The point of this project is to make a program that can not only verify whether or not a given configuration of the Lonely Runner conjecture holds, but also to strengthen our belief in the Lonely Runner conjecture. The program must therefore be able to test a large number of configurations in short time. I have based these on the results found at ``The Computer Language Benchmarks Game''\footnote{The site can be found at http://shootout.alioth.debian.org/u32q/which-programming-languages-are-fastest.php and was last accessed the 7/11 2010}.
\item[Memory footprint:] Memory is slightly less essential than speed, but if the memory footprint is smaller, then we can have more runners in a given configuration. A smaller memory footprint would also seem to indicate that the program will be faster, since less has to be allocated and cleaned up. 
\item[Already implemented data-structures:] I would rather avoid having to implement, debug, and optimise the Priority Queue, and instead relying on a built-in implementation that has proved itself to be free of bugs.
\item[My familiarity with the language:] Familiarity with the language would enable me to get create a bug free implementation of the algorithms faster.
\end{description}

Based on these criteria I will give each programming language a value from 1 to 3, with 3 being the highest, and 1 being the smallest. 

For this project I have considered 3 programming languages: Java, Haskell and C++.

\begin{tabular}{l|c|c|c|c}
        & Speed & Memory & Data-structures & Familiarity \\
\hline
Java    & 2     &  1     &  3              &   3  \\
Haskell & 1     &  2     &  1              &   1  \\
C++     & 3     &  3     &  3              &   2  \\
\end{tabular}

\subsubsection{Conclusion}
Based on this I would argue that the best choice would be C++ on the grounds that it clearly scores the highest, and that it will be easy for me to implement the Geometrical and Numerical algorithms.

\subsection{Problems during implementation}
In this section I will describe some of the major problems I encountered during my implementation of the algorithms, and what I have done to resolve them.

\subsubsection{Rounding errors}
Since the Lonely Runner conjecture in its nature deals with rational numbers, it is clear that this has to be taken into account when implementing it in a real programming language. During my earliest testing phases I quickly discovered that even if the Geometrical algorithm found a valid solution, it would not be valid according to \eqaref{eqa:lonelyRunner}. My solution to this was to make the check entirely discrete. As stated in \eqaref{eqa:lonelyRunner}, we have a Lonely Runner iff
\eqa{
\Vert w_i x\Vert \geq \frac{1}{n+1}
} 

The main problem here being $\Vert w_i x\Vert$ - however, since we can be sure that the solution will always be rational (since all the speeds are integers) we have that x can be written of the form $\frac{p}{q}$ for a given $p \in \N\cup\E{0}$ and $q \in \N$ (since all speeds must be positive and the runners cannot cover negative distance). So now we have

\begin{equation}
\label{compareRewriteOne}
\begin{split}
\Vert w_i * \frac{p}{q}\Vert &\geq \frac{1}{n+1}\\
\min(\frac{w_i * p \bmod q}{q}, 1 - \frac{w_i * p \bmod q}{q}) &\geq \frac{1}{n+1}\\
\min(\frac{w_i * p \bmod q}{q}, 1 - \frac{w_i * p \bmod q}{q}) * (n+1) &\geq 1
\end{split}
\end{equation}

Depending on whether $\frac{w_i * p \mod q}{q}$ or 1 - $\frac{w_i * p \mod q}{q}$ is the smallest value, the equation then becomes

\begin{equation}
\label{compareRewriteTwo}
\begin{split}
\frac{w_i * p \bmod q}{q} * (n+1) &\geq 1\\
(w_i * p \bmod q) * (n+1) &\geq q 
\end{split}
\end{equation}

or 

\begin{equation}
\label{compareRewriteThree}
\begin{split}
1 - \frac{w_i * p \bmod q}{q} * (n+1) &\geq 1\\
(q - (w_i * p \bmod q)) * (n+1) &\geq q
\end{split}
\end{equation}

\subsection{Use of real numbers in the algorithm}
As detailed in Section \ref{eventPoints} I choose to implement the comparison of the Event Points through integer comparison, on the grounds that it is much faster than working with floating numbers, on most hardware. It is clear from \eqaref{eqa:integerTime} that neither side of the expression can be wholly precomputed, as it requires the other Event Point's speed, meaning that this solution trades floating point comparison with several integer multiplications. It is however clear on inspection that avoiding floating point number greatly speeds up the computation.

\subsection{Scalability of implementation}
\label{scale}
An important question that has to be answered is how well the 2 algorithms scale. In this section I will interpret scaling in 2 senses:
\begin{enumerate}
\item How much time it requires to give an answer given a certain input size
\item How large input sizes the algorithm can handle
\end{enumerate}

For the first interpretation it is clear from Section \ref{proof_geo} and \ref{proof_num} that both algorithms have a bad worst case time complexity. These are however unlikely to be relevant during most verification attempts (as we have yet to find a counter example to the Lonely Runner conjecture after more than 2 decades), and therefore not really a good measurement. 

The next question to ask is whether or not the algorithms could be parallelised, to take advantage of the multiple cores we are seeing in today's modern computers. Due to its shared state (the number of runners in the Zone) and the iterative nature of the Geometrical algorithm, I do not believe that the main loop of the algorithm as implemented could be parallelised. However, as discussed in \ref{eventPoints}, there might be a variant where parts of the algorithm could be parallelised. The numerical algorithm on the other hand is ripe for paralleisation, all 4 loops can easily be parralised, having no real shared state (other than, since the algorithm is based on enlightened guesses.

Since the point of this program is to strengthen our trust in Lonely Runner conjecture, the maximum number of runners, and the maximum speed of each runner, that we are able to test is important. I will therefore in the following section analyse the bounds on these values, and discuss how I have improved these.

\subsubsection{Geometrical algorithm}
In the current implementation of the comparison operator is that there is a limit on the number of runners and their speeds. Let us observe  
\eqa{
\label{eqa:firstCompare}
(a + k * (n+1)) * s
}
 from \eqaref{eqa:integerTime}, where all the symbols has the same meaning as before. The possible large numbers are $a$, which in the worst case is $n$, $n$ itself and $s$. No matter which built in number type we choose to represent the above expression with, it is clear there exists values for $n$ and $s$ that will cause an arithmetic overflow.

\hide{
The best I have been able to come up with is rewriting the entire comparison to 
\eqa{
\label{eqa:secondCompare}
s_ia_j - s_ja_i \leq (n+1)(k_is_j - k_js_i)
}
which still requires $s$ to be multiplied with $a$, making the speed of the runner and the local position (which in the worst case is $n$) the bottleneck.
}

The question then becomes what the maximum allowable value becomes by \eqaref{eqa:firstCompare}. Since we are not sure whether or not the Lonely Runner conjecture holds, we must make sure that the expression does not cause arithmetic overflow, even if the algorithm for that configuration of runners will terminate by finding \comFin. We must therefore assume the maximum values for the variables that change - the local position $a$ and k, the number of times the runner has completed a lap. The biggest value of $a$ is $n$ and given a speed $s$, the runner can at most run $s$ rounds before the final runner arrives at the the finishing line. Let $d$ be the maximum positive number that a given number type can represent, then to find the maximum value of $s$ with respect to $n$ and $d$, we have the expression 
\eqa{
(n + s * (n+1)) * s &=& d\\
s^2(n+1) + sn &=& d\\
s^2(n+1) + sn -d &=& 0\\
}
It is clear that this a quadratic equation with $s$ as the unknown - using the standard technique this gets us
\eqa{
s = \frac{-n \pm \sqrt{n^2 + 4(n+1) * d}}{2(n+1)}
}
In practise we are only interested in
\eqa{
s = \lfloor\frac{-n + \sqrt{n^2 + 4(n+1) * d}}{2(n+1)}\rfloor
}
as the other answer would be negative, which has no meaning in this context.

So if $n = 1000$ and $d = 2^{32}-1$ (an unsigned 32 bit integer) then $s = 2070$.

If we try to find $n$ with regards to $d$ and $s$, then we have
\eqa{
s^2n + sn + s^2 &=& d\\
n(s^2 + s) + s^2 &=& d\\
n(s^2 + s) &=& d - s^2\\
n &=& \frac{d - s^2}{s^2 + s}
}
which becomes
\eqa{
n = \lfloor \frac{d - s^2}{s^2 + s} \rfloor
}

So if we let $s = 1000$ and $d = 2^{32}-1$ then $n = 4289$

This is rather constraining, but if we use a 64 bit integer was to be used instead, then combinations such as $n = 10^6$ and $s = 4294964$ would become possible.

Since I in Section \ref{specs} detailed that I wanted the program to run on a normal computer, which yet does not cover 64-bit architectures, I have made the comparisons, and other checks that might cause overflow, use arbitrary-precision arithmetic\footnote{Based on the NTL library, see Section \ref{third_party}} - limiting the possible comparisons only by the memory of the system, at the cost of speed.

\subsubsection{Numerical algorithm}
A problem in the numerical implementation is that no sum $k$ of any 2 speeds must exceed the maximum value of the numerical type used (say, $2^{32} - 1$ for 32 bit integers). Even if this is far better than the Geometrical algorithm, I have implemented arbitrary-precision arithmetic, in order to make the algorithms equal in this regard.

\subsection{Third party libraries used}
\label{third_party}
\begin{description}
\item[NTL arbitrary-precision arithmetic (v. 5.5.2 ):] A software library to allow for arbitrary precision sized integers in both algorithms. The library can be found at http://www.shoup.net/ntl/doc/tour.html (last visited 10/11/2010).
\item[The GNU Multiple Precision Arithmetic Library (GMP) (v. 5.0.1):] The NTL library has been compiled towards the GMP to increase its efficiency. The library can be found at http://gmplib.org/ (last visited 10/11/2010).  
\item[json-c (v. 0.9):] A C++ library used for importing and exporting data to/from the JSON file format. The library can be found at http://oss.metaparadigm.com/json-c/ (last visited 10/11/2010).
\end{description}

\section{Test Results}
\label{test}
\todo{make tests and describe them}

\subsection{Introduction}
In this section I will introduce how I am going to test the algorithms, which parameters I am going to test them on, which software and hardware setup I have tested them on, description of the test themselves, and finally the results.

\subsection{How to test the Algorithms}
Since a configuration of the Lonely Runner conjecture requires $n$ runners and each runner needs a distinct speed, it is clear that the parameters we must test the algorithms for are the number of runners, and their speed. Since the speeds have to be different, I will use the maximum speed to be the defining factor.

To do this test in practice I have choosen to do it in increments, with the number of runners being the deciding factor. If we have to test for $n$ runners, then we let the first configuration to test contain the first $n$ numbers, the second contain from ($n+1$)'th number to the ($2n+1$)'th number and so on.

For example if we have to test 10 runners for values up to 100 with sequential numbers, then the first array to be tested will be [2, 3, 4, $\ldots$, 9, 10, 11], the second array to be tested [12, 13, 14, $\ldots$, 19, 20, 21] and so on.

Since the Lonely Runner conjecture is about the relation between numbers, it would be interesting to test the algorithms in the cases where the numbers are pairwise prime. To do this I have simply calculated the prime numbers up to the maximum value, and used them as the speed.

As a stress test I have also tried calculating the sequential numbers for the 

\subsection{Hardware and software setup}
As mentioned in the implementation part I have implemented the software in C++ using g++ 4.4.3 compiler for Ubuntu. The test was run on a [] \todo{name computer type, hardware, os and so forth once final tests have been made}.

\subsection{Test made} 

I will test both alorithms with 10, 50, 100, 500, 1000 and 2000 runners with a max speed up to 5000 - both for sequential,and prime number speed. For a given configuration the number runners will of course be their minimum speed

\section{Results}
\label{results}
\todo{make results}
In this section I will show and discuss the results I have had running the program against various data-sets. As specified in the last section, I have made 3 different sets of data: Prime Numbers, Sequential Numbers and Randomly generated numbers. Since each data-set follows a similar pattern, I have chosen to talk about the most important of these, instead of going through each of the data-sets and repeating the commentary. 


For brevity's sake I have only included the most important sets of data for reasons of brevity. Complete listings of the data-sets in table form can in \ref{} and will be made available in machine-readable form upon request.


\subsection{Small number of runners}
\graph{test_Sequential-Random_50_runners_Geo}{}

\graph{test_Sequential-Random_100_runners_Geo}{}

\graph{test_Sequential-Random_500_runners_Geo}{}

[prime], [seq], [ran-sorted]:

It is clear from these graphs, that as long as there are less than 500 runners, then the Numerical algorithm is faster than the Geometrical. However, around 500 runners something happens and for the first few speeds the Numerical algorithm takes an unreasonable time to complete (more than 1 second). In depth study of the execution of the Numerical algorithm has shown that in these cases, finding a solution is not done for the first few combinations of $k$-values. I believe this to be an effect of the fact that the Numerical solution does not systematically try to find a solution, but rather tests possible times. Thus the Numerical algorithm is more of a gamble, but the numbers clearly show that most of the time, it's better than the Geometrical.

Since the Numerical algorithm is clearly input dependent, I have tried to randomize the input, in hopes that this would yield better results. However, as can seen below, this does not make the time spikes go away, but merely moves them around.

[randomized values]




 

 that the Numerical algorithm is generally faster than the Geometrical algorithm for configurations that contains speeds under 20000. However it 

It is also clear, that while the Geometrical algorithm is clearly slower, it still only need microseconds to find a solution. 

\subsection{Larger numbers}
 
For these cases it is clear that the values are 

\subsection{Conclusion}


% [Conclusions are very important. Do not expect that the reader remembers everything you told him/her.
% Having stated the definitions, you can now be more specific that  in the introduction]
% * Overview what this work was about.
% * Main results and contributions
% * Comments on importance or
% * Tips for practical use [how your results or experience can help someone in practice or
%     another researcher to use your simulator or avoid pitfalls]
% * Future work. Reinforce the importance of work, but avoid giving out your ideas].

\section{Conclusion}
\label{conclusion}

\subsection{Summery}

I have in this project:
\begin{itemize}
\item Implemented and tested the Numerical method found in \cite{invis}
\item Implemented and tested a Geometrical method of my own devising.
\item Argued that if there exist a time that makes \eqaref{eqa:lonelyRunner} hold, then it will be rational.
\end{itemize}

\subsection{Results}


\subsection{Future work}
\begin{description}
\item[Improved Geometrical method:] It would be interesting to see whether the Geometrical algorithm could be made faster. One way could be only investigate the intervals where the slowest runner is in the Zone. This is based on the fact that we only have a valid solution if all the runners are in the Zone, which more specifically requires the slowest runner to be in the Zone. This could quite possibly dramatically reduce the required run-time, as since all creations, insertions and removals of Event Points when the runner is outside the Zone would never be made. While the modular nature of the Lonely Runner conjecture might make this seem like a futile task, it is important to remember that it is possible calculate the times where a runner enters the zone, and thus 

Let SlowZone be the interval where the slowest runner is in the zone, then a pair of Event Points should be inserted on the heap iff at least one of \comStart and \comFin is in SlowZone. This would however separate the algorithm into 2 parts: The first part which finds the first Event Points that could be interesting, and the second part which is the same as before, only with a check that does not insert pairs of Event Points that are entirely outside the SlowZone.

\item[Parallelism:] It is clear that the Numerical algorithm can take advantage of multiple threads. However, given the results, it can be debated whether this would significantly improve its run-time, given that it either finds the solution right of the bat, or has to run for several hours (in which case one might wonder whether it might be better to make a hybrid algorithm that runs both, that stops them both when a solution has been found). A better place to utilize parallelism would perhaps be in the Range test, distributing the intervals that each CPU (or indeed, computer) should check. This is especially true, since the memory peak load of the range test is very small.

\item[Verification for many different number of runners and speeds:] Since I in this project has made a verifier for whether a given configuration holds for \eqaref{eqa:lonelyRunner} or not, it seems that the next logical step would be to set it to verify the Lonely Runner conjecture for a large verity of runner configurations. 
\end{description}

\bibliography{LonelyRunner}{}
\bibliographystyle{plain}
\appendix
\section{Tables}
\label{data_tables}
\includeTable{testPrimes10runners}


\hide{
\section{Code}
}
\end{document}
